\section{Choosing the numerical core before asking an agent to choose}
\label{exec14:l4-alternatives}

The proposed protected core still leaves a substantive choice. Constraint
search, score optimization, functional equations and changes of environment
can all return a graph, but they use different information to distinguish
directions. An assistant should explain the reason for choosing among them.
Selecting the package that happens to accept the input table is insufficient.
For a branch-pricing problem, the first questions are whether the table records
independent units or a dependent panel, whether omitted common causes are
plausible, whether the action and outcome have a reliable order, and whether
the proposed structural restrictions are defensible for those variables.
Only then is it useful to compare runtime and interface convenience.

PC uses conditional-independence information with a causally sufficient DAG
interpretation under its assumptions. Its partially directed output preserves
orientations that the available independence information does not determine.
FCI permits latent and selection variables under a different model class and
returns PAG information. RFCI trades some adjacency information for a less
expensive search. The differences matter when an analyst asks whether a
borrower's private assessment jointly influences an offer and subsequent
recovery. Selecting PC after simply deleting that unrecorded assessment from
the variable list does not remove the common cause. Selecting FCI does not
make the effect point identified. Its output may correctly retain precisely
the ambiguity that prevents a unique intervention calculation
\citep{exec13_latentr_colombo2012,exec14_java_latent_sources}.

A score search instead compares how well candidate structures represent the
data after accounting for complexity or prior preferences. A selected DAG
can be an effective computational summary without being the only causal
explanation. The score, its statistical interpretation, the search space and
the optimization outcome are separate choices. If two structures have the
same observable implications under the supplied information, an arbitrary
tie break does not become a discovered direction. A language prior can break
the tie, but the report must then attribute the direction to that added
restriction or preference. For the proposed bank tool, a ranked candidate
list is useful if its conditional meaning remains visible; an unsupported
claim that the first entry is the true mechanism is not.

NOTEARS illustrates both the appeal and the limit of a continuous formulation.
For data matrix $X$ with $n$ rows and $d$ columns, the inspected linear
formulation uses a coefficient matrix $W$ and the objective
\[
 F(W)=\frac{1}{2n}\lVert X-XW\rVert_F^2+\lambda\lVert W\rVert_1,
 \qquad h(W)=\operatorname{tr}\exp(W\circ W)-d=0.
\]
Here $W_{ij}$ represents a coefficient from variable $i$ to variable $j$,
$\circ$ denotes entrywise multiplication, and the exponential is a matrix
exponential. This convention differs from a target-row coefficient matrix.
The acyclicity characterization is an inherited result, not an identifying
assumption that language can supply. If $S=W\circ W$, every entry of $S$
is nonnegative and each $\operatorname{tr}(S^k)$ sums weighted closed walks.
The expansion $\operatorname{tr}\exp(S)=d+\sum_{k\geq1}
\operatorname{tr}(S^k)/k!$ is equal to $d$ exactly when the graph has no
directed cycle. Squaring the entries prevents positive and negative edge
weights from cancelling a cycle in this argument
\citep{notears2018}.

That theorem separates a structural requirement from the inferential question.
A small computed $h(W)$ is a numerical diagnostic of one constraint. It is
not a probability that the graph is correct. The original paper expressly
distinguishes stationary-point optimization from global minimization, and
its statistical interpretation refers to assumptions studied separately.
The inspected reference implementation also has an $\ell_1$ penalty,
augmented-Lagrangian stopping rules and a final weight threshold. Each can
change the output. An agent that raises the threshold until a graph looks
sparse has changed the analysis; it has not improved a calibrated confidence
level. The exact returned graph needs a structural check after thresholding,
and the fit needs an appropriate scientific interpretation
\citep{notears2018,exec13_pydisc_notears}.

An apparent advantage of one discovery method can also depend on units.
Multiplying a deposit balance column by a thousand changes an unstandardized
squared-error objective and the relative meaning of a common coefficient
penalty. A nonlinear transformation changes the functional model still more.
The transformation, back-transformation and intervention scale should be
part of the proposed specification. A rank-based numerical method, a linear
equation method and a language judgment need not have the same invariances.
The comparison should expose this difference instead of assuming that an
identical CSV file makes all models scientifically matched.

LiNGAM obtains directional information from a different restriction: a linear
acyclic system with mutually independent, non-Gaussian disturbances and no
unobserved confounders. Its original model assumes repeated independent
vectors generated with the same equations and disturbance distributions.
The inspected original construction uses independent component analysis;
the currently inspected DirectLiNGAM code is a separate algorithmic route.
The source explicitly notes that this model need not impose faithfulness.
Consequently, a failure of a faithfulness-based argument does not by itself
refute LiNGAM. Conversely, apparent skewness of observed loan balances does
not establish the non-Gaussianity and independence of the structural
disturbances in the chosen equations. These are requirements on the model,
not labels that can be inferred from column names
\citep{exec13_pydisc_shimizu2006,exec13_pydisc_lingam}.

The Gaussian equivalence example earlier in this chapter intentionally lies
outside that non-Gaussian restriction. It cannot be used as an impossibility
proof for every functional discovery model. This distinction is productive:
one can escape an observational equivalence by supplying additional
information or by restricting the admissible structural models. The proposal
should state which of these it is doing and what evidence supports the
restriction. The language component can retrieve a scientific argument for
linear mechanisms, question an independence assumption or compare an
alternative model. It cannot verify disturbance independence merely by
producing a coherent explanation.

For time series, the underlying sampling and causal model changes again.
The inspected PCMCI+ result requires its stated sufficiency, stationarity,
faithfulness, lag coverage and test consistency conditions; selection and
measurement error are outside that theorem. Contemporaneous directions can
remain unresolved. A calendar panel containing common shocks and repeated
accounts cannot inherit the result solely because the records have a date
column. The assistant should identify the sampling unit and candidate lag
family, explain which common causes are represented, and retain uncertainty
about contemporaneous effects. Declaring a lag of one because the prompt is
shorter is a substantive restriction, not a harmless interface choice
\citep{runge2020}.

These comparisons imply a staged software decision. Retain method-specific
profiles for PC, FCI/RFCI, functional equations, continuous score search and
temporal discovery. Each profile binds the statistical assumptions, graph
meaning, matrix convention, search bounds and failure behavior. Reuse a
reviewed implementation when its complete path can be made reproducible;
consider porting only when a measured deployment or numerical requirement
justifies the maintenance burden. A common result schema can preserve their
differences. It should not force every method to emit a fully oriented DAG or
an effect estimate when that is stronger than the method supports.

\section{Faithfulness can fail even with complete measurements}
\label{exec14:l4-faithfulness}

A second exact project example separates missing variables from cancellation.
Let $A,E_M,E_Y$ be independent centered Gaussian variables with unit variance,
and define
\[
 M=A+E_M,\qquad Y=-A+M+E_Y=E_M+E_Y.
\]
The structural graph contains $A\to M$, $A\to Y$ and $M\to Y$. All three
variables are observed, all disturbance variances are positive, and the
joint distribution has a nonsingular covariance matrix. Nevertheless
$A\perp Y$: the direct and mediated paths cancel in the marginal law.
The covariance matrix, in the order $(A,M,Y)$, is
\[
 \Sigma=\begin{pmatrix}1&1&0\\1&2&1\\0&1&2\end{pmatrix},
 \qquad
 \operatorname{Cov}(A,Y\mid M)=0-\frac{1\cdot1}{2}=-\frac12.
\]
Thus independence at the empty conditioning set and dependence after
conditioning can arise here without the structural collider explanation
that a faithful model would use. This is not a numerical roundoff example.
The zero is an exact property of the stipulated population law.

There is indeed a different structural model with that faithful collider
explanation. Take mutually independent $A\sim N(0,1)$, $Y\sim N(0,2)$
and $\eta\sim N(0,1/2)$, and set
\[
 M=A+\tfrac12Y+\eta.
\]
It produces exactly the same $\Sigma$ and therefore the same centered
Gaussian distribution. In the triangle model, changing $M$ while leaving
$A$'s equation intact changes the mean of $Y$ with derivative one. In the
collider model, intervening on $M$ leaves $Y$ unchanged. The target effects
are different despite complete observed-variable measurement in both models.
Invoking faithfulness selects a model class that excludes the triangle; it
does not make the observable covariance itself reveal which structural world
generated the records \citep{exec14_discovery_design_oracles}.

For a banking interpretation, $A$ may represent standardized offer intensity,
$M$ a subsequent account-balance response and $Y$ net revenue. A direct cost
and an indirect revenue channel can offset each other. The equations are
stipulated and should not be taken as an empirical model of that mechanism.
They explain why an observed zero association can have several meanings.
A reliable chronology establishing $M$ before $Y$ would rule out the second
model as an ordinary acyclic causal explanation. That chronology is additional
information. With only a coarse accounting period, an agent must not assume
that the order of columns supplies it.

Cancellation is also a finite-precision concern. Replace the direct
coefficient by $-1+\delta$. Then
\[
 Y=\delta A+E_M+E_Y,\qquad
 \operatorname{Cov}(A,Y)=\delta,
 \qquad \operatorname{Cov}(A,Y\mid M)=\frac{\delta-1}{2}.
\]
The marginal signal can be arbitrarily small while the conditional signal
remains appreciable. A consistency statement for a fixed nonzero signal does
not guarantee useful uniform recovery as $\delta$ approaches zero. This
explains why the prototype should vary cancellation strength and retain
weak-signal failures. Increasing prompt detail does not increase the numerical
separation between these distributions.

The appropriate remedy depends on the claim. A verified time order can exclude
some candidate graphs. An intervention on the consequential variable can
supply new information. A functional model can use restrictions not present
in conditional-independence information, provided those restrictions are
defensible. A conservative report can retain multiple mechanisms and compare
their decision implications. None of these remedies is equivalent to choosing
an orientation because it sounds natural in a bank narrative. The exact
example supports a test of this distinction without paying for live calls or
obtaining bank data.

\section{A search budget changes which separations can be found}
\label{exec14:l4-depth-example}

The conditioning-set budget has a similarly concrete scientific meaning.
Consider mutually independent unit-variance centered Gaussian disturbances
and the structural equations
\[
 A=E_A,\qquad M_1=A+E_1,\qquad M_2=A+E_2,
 \qquad Y=M_1+M_2+E_Y.
\]
There are two parallel paths from $A$ to $Y$ and no direct edge. Their total
intervention derivative is two. The observed covariance matrix is
\[
 \Sigma_{A,M_1,M_2,Y}=
 \begin{pmatrix}
 1&1&1&2\\1&2&1&3\\1&1&2&3\\2&3&3&7
 \end{pmatrix}.
\]
The Gaussian conditional covariance calculation gives
\[
 \operatorname{Cov}(A,Y)=2,\quad
 \operatorname{Cov}(A,Y\mid M_1)
 =\operatorname{Cov}(A,Y\mid M_2)=\tfrac12,
 \quad \operatorname{Cov}(A,Y\mid M_1,M_2)=0.
\]
For the last equality, the mediator covariance inverse is
$\frac13\left(\begin{smallmatrix}2&-1\\-1&2\end{smallmatrix}\right)$;
multiplying the cross-covariances subtracts exactly two. A search limited to
conditioning sets of size at most one cannot find the separator for this
pair, even with an exact oracle and an unlimited sample. A reported extra
adjacency is therefore a possible consequence of the search restriction,
rather than evidence for a direct offer-to-revenue mechanism
\citep{exec14_discovery_design_oracles}.

This example also distinguishes a discovery query from an adjustment rule.
Conditioning on the mediators helps discover that a direct adjacency is
unnecessary. Adjusting them away when estimating the total effect of $A$
would remove the very paths that define that effect. An agent should not copy
the variables from a separating set into an outcome-regression adjustment
list without reviewing the target and graph. The same named set can be
appropriate for one computational question and inappropriate for another.

The proposed comparison should therefore bind the maximum conditioning order,
the admissible variable sets, the CI test and its missing-result behavior.
Record the number of attempted queries as well as completed ones and preserve
which search bounds were active. A faster configuration with smaller search
depth is a useful alternative when the user accepts its scope, but it is not
automatically a faster implementation of the same inferential procedure.
For banking data, higher conditioning order can also reduce usable observations
and aggravate singularity or mixed-type test limitations. Relaxing one
computational constraint can create another statistical problem.

This small model is a conformance reference for a future integration. It can
check that the adapter preserves query variables, reports the configured
depth and distinguishes a population zero from an unavailable test. The
calculation does not claim that any inspected package has been run or that
the entire search will visit a particular query in a particular order.
Whole-engine qualification still needs a pinned executable path and an
observed trace. The source-inspected Tetrad branch discussed earlier and this
mathematical example answer complementary questions.

\section{Several environments need a causal interpretation of their changes}
\label{exec14:l4-environments}

An assistant may propose using changes across branches, countries or policy
periods to orient a relation. Invariant prediction provides a substantive
alternative to single-environment graph search. The inspected formulation
assumes that, for a suitable predictor subset, the response mechanism and
independent error law remain invariant across environments. In the linear
case, the same coefficients and residual distribution apply in each setting.
Under its stated test and candidate-set assumptions, the intersection of
nonrejected subsets controls false inclusions in that invariant subset; it
does not promise recovery of every cause. Connecting the invariant subset to
causal parents requires the corresponding structural assumptions. In
particular, arbitrary changes to the response or its hidden causes are not
licensed by the presence of an environment label
\citep{petersicp2016}.

The distinction can be exhibited with another exact project law. Index
environments by a known number $m$. Let independent variables satisfy
$U\sim N(m,1)$, $E_A\sim N(0,1)$ and $E_Y\sim N(0,1/2)$, and write
\[
 A=U+E_A,\qquad Y=U-\frac m2+E_Y.
\]
Here $A$ has no causal effect on $Y$. Yet in every environment
\[
 E(A)=m,\quad E(Y)=m/2,\quad
 \operatorname{Var}(A)=2,\quad \operatorname{Cov}(A,Y)=1,
 \quad\operatorname{Var}(Y)=3/2.
\]
Gaussian conditioning gives $Y\mid A,m\sim N(A/2,1)$, independent of $m$.
A stable conditional predictor has emerged from changing a hidden common
cause and compensating the outcome intercept. The different model
$A\sim N(m,2)$, $Y=A/2+\eta$ with independent $\eta\sim N(0,1)$ has
the same observed law in every environment and an intervention derivative
of $1/2$. Thus the environment-indexed observations do not distinguish these
two causal explanations \citep{exec14_discovery_design_oracles}.

This construction does not contradict the invariant-prediction theorem. It
violates the structural interpretation needed to treat the stable observed
regression as the unchanged causal response mechanism: the hidden cause and
the outcome equation change together. It demonstrates why a test of one
observed conditional distribution cannot by itself certify the provenance of
the environment changes. An apparently successful comparison across many
branches can still be uninformative if their changes share this structure.
The number of environments is not a count of independent randomized
interventions on the proposed action.

For a deposit application, a market stress period may change demand, the
bank's pricing rule and measurement practices at the same time. A dated
policy document can establish that pricing changed, but it may not establish
that all other relevant mechanisms remained fixed. A useful language task is
to construct an environment dossier: what changed, when it became effective,
which variables it could affect, and what the sources leave unknown. If the
dossier supports a narrower intervention family, a numerical method can be
assessed under that family. Otherwise the environment comparison remains a
predictive or sensitivity analysis with explicit limits.

The public prototype can test this reasoning affordably. Generate environments
that alter only an admissible predictor mechanism and compare them with the
compensating-change construction above. Keep the target and permitted
documents fixed within each regime, and hide the generator from the assistant.
A successful workflow should distinguish failure to reject invariance from
evidence that the structural assumptions were supplied. A correct refusal in
the second regime and a useful result in the first are different outcomes
that should both appear in the evaluation. This is a proposed study design;
the exact algebra supplies reference cases without supplying performance
estimates for any model.

\section{Allowing mistakes in extracted restrictions without hiding them}
\label{exec14:l4-restriction-errors}

Hard restrictions become especially fragile when an agent extracts many of
them. Let $\mathcal G_D$ be the graph class retained from numerical evidence,
and let $R_1,\ldots,R_J$ be proposed predicates on a graph. The ordinary
intersection retains only graphs satisfying every predicate. A more explicit
sensitivity construction allows at most $k$ predicates to be wrong:
\[
 \mathcal G_{D,k}=\left\{G\in\mathcal G_D:
       \sum_{j=1}^J\{1-R_j(G)\}\leq k\right\}.
\]
This is a proposed deterministic construction, not an implemented general
graph solver. If the true graph lies in $\mathcal G_D$ and violates at most
$k$ proposed restrictions, it lies in $\mathcal G_{D,k}$. Enlarging $k$ cannot
shrink the retained class. Mapping that class into effect sets makes the
decision's dependence on the trusted restrictions visible.

For an exact finite illustration, label four otherwise admissible candidate
models $G_{ab}$ by two binary properties $a,b$. Suppose the extracted claims
assert $a=1$, $b=1$, and that exactly one property is one. The joint
intersection is empty. Allowing one erroneous claim retains $G_{01}$,
$G_{10}$ and $G_{11}$, since each violates exactly one claim; $G_{00}$
violates two and is excluded. A nonempty relaxed class has therefore not
proved that the true model was retained. If the truth is $G_{00}$, the
one-error sensitivity analysis still misses it. The remedy requires either
more trusted evidence or a larger declared error allowance, with the
consequent loss of precision \citep{exec14_discovery_design_oracles}.

A probabilistic guarantee requires actual error control in addition to that
logical construction. Write $N=\sum_j\mathbf1\{R_j(G_\star)=0\}$ for
the number of wrong restrictions at the true graph. If a justified evaluation
or observation law supplies marginal bounds
$\Pr\{R_j(G_\star)=0\}\leq\delta_j$, then
$E(N)\leq\sum_j\delta_j$. Since
$\mathbf1\{N>k\}\leq N/(k+1)$, it follows that
\[
 \Pr\{G_\star\notin\mathcal G_{D,k}\}
 \leq \alpha_D+\frac{\sum_j\delta_j}{k+1},
 \qquad
 \Pr\{G_\star\notin\mathcal G_D\}\leq\alpha_D,
\]
with the right side capped at one. Independence between the restrictions,
or between text and numerical evidence, is unnecessary for this inequality.
Its premises must concern the actual joint workflow, including adaptive
extraction. A model's confidence statement or vote frequency is not a
supplied $\delta_j$. The result is an elementary project deduction exposing
what a reliable probabilistic claim would need.

For three genuine marginal bounds of $.05$ each, the restriction component
is at most $.15$ when $k=0$ and $.075$ when $k=1$. These values are conditional
illustrations; no such language-model error bounds have been established in
this proposal. A larger retained class can make the effect or decision
unresolved even if its retention bound improves. This separates the value of
error tolerance from a promise of useful precision.

The unit of an error allowance needs substantive care. Three copied passages
can express the same mistaken policy statement. Treating them as three
independent confirmations is unjustified, while counting them as three
different restrictions can also distort an at-most-$k$ analysis. The proposed
system should retain a mapping from predicates to source assertions and
source families. An analyst can then define sensitivity to a failed source
family rather than to a number of duplicated sentences. A newly discovered
contradiction should create a revised analysis, preserving the earlier
restrictions and showing which effects or decisions changed.

\section{An intervention's value can be calculated before conducting it}
\label{exec14:l4-finite-information}

The perfect-information calculation above is an upper bound. A prospective
experiment normally supplies noisy information, so a prototype should also
evaluate the attainable value under a declared experiment. Return to the
chapter's two binary worlds. In the forward world, $A$ is a fair binary action
and $Y$ equals $A$ with probability $.9$. In the reverse world, $Y$ is fair
and $A$ equals $Y$ with probability $.9$. Both have the same observational
table. Under a controlled action $A=1$, however, the outcome success
probabilities are $.9$ and $.5$, respectively. This difference comes from
replacing the action equation, not from resampling the old observational
table.

Assume, solely for this planning calculation, that independently intervened
units can be drawn from the same structural world, that outcomes are observed
without loss, and that neither interference nor experimentation changes the
mechanism. If $K$ successes are observed in $n$ probes, let
\[
 f_H(k)=\binom nk(.9)^k(.1)^{n-k},\qquad
 f_L(k)=\binom nk(.5)^n.
\]
For an elicited prior probability $q$ on the forward world, Bayes' rule gives
\[
 q_k=\frac{qf_H(k)}{qf_H(k)+(1-q)f_L(k)}.
\]
The common binomial coefficient cancels in the likelihood ratio, so the
posterior odds equal the prior odds multiplied by
$1.8^k\, .2^{n-k}$. Unlike a prompt-vote fraction, this update has an
explicit observation law and an explicit relationship to the two worlds.
Its validity remains conditional on those two worlds and the stipulated
intervention protocol.

Connect it to the earlier offer decision by assigning gross incremental value
ten times the intervention effect and an offer cost of three. The effects
are $.8$ and zero, giving net offer values five and minus three; no offer
has value zero. For a future comparable account, the optimal action after
observing $k$ is to offer when $q_k>3/8$. Before the probe results are known,
the expected optimal value is
\begin{equation}
 V_n(q)=\sum_{k=0}^n
       \max\{0,\,5qf_H(k)-3(1-q)f_L(k)\}.
 \label{exec14:l4-sample-information-value}
\end{equation}
To obtain this expression, multiply the conditional optimum
$\max(0,8q_k-3)$ by the predictive mass
$qf_H(k)+(1-q)f_L(k)$ and simplify. The resulting sum includes every possible
outcome, including evidence that leads to no offer. There is no favorable
outcome selection.

At $q=1/2$, the no-information value is one. With one probe, a success leads
to an offer and a failure leads to no offer; the expected value becomes
$3/2$, so the expected value of sample information is $1/2$. With two probes,
only two successes justify an offer. The expected value is $33/20$, giving
sample-information value $13/20$. Exact enumeration for larger $n$ appears
in Table~\ref{exec14:l4-evsi-table}. These are arithmetic consequences of a
stipulated finite model, not observed experimental results or measured bank
benefits \citep{exec14_discovery_design_oracles}.

\begin{table}[htbp]
\centering
\caption{Exact information value under the stipulated two-world intervention law.}
\label{exec14:l4-evsi-table}
\begin{tabular}{rrr}
\toprule
Probes $n$ & Success count triggering an offer & Value of sample information\\
\midrule
0 & $0$ or more & $0.000000$\\
1 & $1$ or more & $0.500000$\\
2 & $2$ or more & $0.650000$\\
5 & $4$ or more & $1.015100$\\
10 & $8$ or more & $1.242492$\\
20 & $15$ or more & $1.440825$\\
50 & $37$ or more & $1.498585$\\
\bottomrule
\end{tabular}
\end{table}


The value is bounded above by the perfect-information value of $3/2$ per
future account. It cannot decrease when another free, comparable probe is
made available, because a decision rule can ignore that probe. This is a
property of the decision problem and information nesting; it does not imply
that a paid study has nonnegative net value. If $N_f$ future accounts share
the relevant world and reward model, the gross planning value is
$N_f\{V_n(q)-V_0(q)\}$. The complete study cost, including probe treatment,
observation, review, delay and any required harm constraints, must be assessed
against that total. Multiplying by future accounts assumes that the same
learned information transfers to those accounts; heterogeneous worlds would
require a different calculation.

This calculation supplies a clear reason to prefer a generated prototype now.
It can verify the likelihood update, target preservation, action threshold
and complete outcome accounting. It can introduce dropout, shared shocks or
worlds outside the two-model family and test whether the workflow reports
the changed assumptions. An actual bank experiment is justified only when
its feasible action, observation law, transfer population and cost could
resolve a consequential uncertainty. The present calculation cannot supply
those inputs and does not authorize or perform an experiment.

A retrieval-only query has zero value for discriminating these two worlds
if its response law is identical in both and it conveys no new applicable
restriction. Its likelihood ratio is one and the posterior remains $q$.
Retrieval may still improve variable definitions, reduce clerical effort or
reveal a genuine distinguishing fact. Those are separate benefits to model
and measure. Repeating a prompt until it favors the preferred graph does
not reproduce the experiment analyzed in Equation~\eqref{exec14:l4-sample-information-value}.

The comparison also explains why full graph recovery is sometimes unnecessary.
When every retained model recommends the same feasible action over the
relevant cost range, further orientation may have little decision value.
When two models differ on a consequential action, the useful query is the
one that separates them under an admissible information mechanism. An agent
can help formulate that query and show which assumptions drive its value.
The proposal's initial objective is this reviewable connection from discovery
to a decision, with numerical and semantic limits preserved.
\label{exec14:l4-completion-end}
